% from ega2-1-fr.tex
\oldpage[II]{37}

quement à $\Gamma(D_+(f),\widetilde{M})$, on obtient ainsi un
homomorphisme de groupes abéliens
\[
  \alpha_0^f:M_0\longrightarrow\Gamma(D_+(f),\widetilde{M}).
\]
Il est clair que, pour tout $g\in S_e$ ($e>0$), le diagramme
\[
  \xymatrix{
    & \Gamma(D_+(f),\widetilde{M})\ar[dd]\\
    M_0\ar[ur]^{\alpha_0^f}\ar[dr]_{\alpha_0^{fg}}\\
    & \Gamma(D_+(fg),\widetilde{M})
  }
\]
est commutatif~; cela signifie que pour tout $x\in M_0$, les sections
$\alpha_0^f(x)$ et $\alpha_0^g(x)$ de $\widetilde{M}$ coïncident dans
$D_+(f)\cap D_+(g)$, et par suite il existe une section unique
$\alpha_0(x)\in\Gamma(X,\widetilde{M})$ dont la restriction à chaque
$D_+(f)$ est $\alpha_0^f(x)$. On a ainsi défini (sans utiliser l'hypothèse
que $S$ est engendrée par $S_1$) un homomorphisme de groupes abéliens
\[
\label{II.2.6.2.1-fr}
  \alpha_0:M_0\longrightarrow\Gamma(X,\widetilde{M}).
\tag{2.6.2.1}
\]

Appliquant ce résultat au $S$-module gradué $M(n)$ (pour tout
$n\in\mathbf{Z}$), on obtient pour chaque $n\in\mathbf{Z}$ un homomorphisme
de groupes abéliens
\[
\label{II.2.6.2.2-fr}
  \alpha_n:M_n=(M(n))_0\longrightarrow\Gamma(X,\widetilde{M}(n))
\tag{2.6.2.2}
\]
(compte tenu de (\hyperref[II.2.5.15-fr]{2.5.15}))~; d'où un homomorphisme
fonctoriel (de degré $0$) de groupes abéliens gradués
\[
\label{II.2.6.2.3-fr}
  \alpha:M\longrightarrow\Gamma_*(\widetilde{M})
\tag{2.6.2.3}
\]
(aussi noté $\alpha_M$) qui, dans chaque $M_n$, coïncide avec $\alpha_n$.

Si on prend en particulier $M=S$, on vérifie aussitôt (compte tenu de la
définition de la multiplication dans $\Gamma_*(\mathcal{O}_X)$
(\hyperref[0.5.4.6-fr]{0, 5.4.6})) que
$\alpha:S\to\Gamma_*(\mathcal{O}_X)$ est un homomorphisme d'anneaux
gradués, et que pour tout $S$-module gradué $M$,
(\hyperref[II.2.6.2.3-fr]{2.6.2.3}) est un di-homomorphisme de modules
gradués.
\end{env}

\begin{proposition}[2.6.3]
\label{II.2.6.3-fr}
Pour tout $f\in S_d$ ($d>0$), $D_+(f)$ est identique à l'ensemble des
$\mathfrak{p}\in X$ où la section $\alpha_d(f)$ de $\mathcal{O}_X(d)$ ne
s'annule pas (\hyperref[0.5.5.2-fr]{0, 5.5.2}).
\end{proposition}

\begin{proof}
Comme $X=\bigcup_{g\in S_1}D_+(g)$ par hypothèse, il suffit de vérifier que
pour tout $g\in S_1$, l'ensemble des $\mathfrak{p}\in D_+(g)$ où
$\alpha_d(f)$ ne s'annule pas est identique à $D_+(fg)$. Or, la restriction
de $\alpha_d(f)$ à $D_+(g)$ est par définition la section correspondant à
l'élément $f/1$ de $(S(d))_{(g)}$~; par l'isomorphisme canonique
$(S(d))_{(g)}\xrightarrow{\sim}S_{(g)}$
(\hyperref[II.2.5.7-fr]{2.5.7}), cette section de $\mathcal{O}_X(d)$
au-dessus de $D_+(g)$ s'identifie à la section de $\mathcal{O}_X$ au-dessus
de $D_+(g)$ qui correspond à l'élément $f/g^d$ de $S_{(g)}$~; dire que cette
section s'annule en $\mathfrak{p}\in D_+(g)$ signifie que
$f/g^d\in\mathfrak{q}$, où $\mathfrak{q}$ est l'idéal premier de $S_{(g)}$
correspondant à $\mathfrak{p}$ (\hyperref[II.2.3.6-fr]{2.3.6})~; par
définition cela veut dire que $f\in\mathfrak{p}$, d'où la proposition.
\end{proof}

\begin{env}[2.6.4]
\label{II.2.6.4-fr}
Soit maintenant $\mathcal{F}$ un $\mathcal{O}_X$ module, et posons
$M=\Gamma_*(\mathcal{F})$~; en vertu de l'existence de l'homomorphisme
d'anneaux gradués $\alpha:S\to\Gamma_*(\mathcal{O}_X)$, $M$ peut être
considéré comme un $S$-module gradué. Pour tout $f\in S_d$ ($d>0$), il
résulte de (\hyperref[II.2.6.3-fr]{2.6.3}) que la restriction à $D_+(f)$ de
la section $\alpha_d(f)$ de $\mathcal{O}_X(d)$ est inversible~; il en est
donc de même de la restriction à $D_+(f)$ de la section $\alpha_d(f^n)$ de
$\mathcal{O}_X(nd)$, pour tout $n>0$. Soit alors
$z\in M_{nd}=\Gamma(X,\mathcal{F}(nd))$ ($n>0$)~; s'il existe un entier
$k\geq0$ tel que la restriction à $D_+(f)$ de $f^kz$, c'est-à-dire la
